\chapter{RAS / ORMAS model-space input (menu 28)}
\label{ch:menu28}
\menulabel{28}

\section{Purpose}

The generalized active space (GAS) concept --- several orbital sub-spaces
with occupation constraints --- enters ORCA as an \emph{incomplete model
space}: the \code{refs} sub-block of \code{\%casscf} takes a partition mask
instead of the complete active space.  ORCA spells the two documented masks
\code{RAS(Nel: NRAS1 MaxHoles / NRAS2 / NRAS3 MaxParticles)} and
\code{ORMAS(nel: m1 min1 max1, m2 min2 max2, ...)} (up to 25 sub-spaces).
The menu generates the input for either mask on either route --- the
orbital-optimized \code{\%casscf} route (RASSCF / ORMAS-SCF) or the
CI-only \code{\%rasci} module; the engine run stays with the user.

\section{What you need}

A structure (XYZ) and the partition mask.  The mask is the single source of
the active-space numbers: the RAS mask's orbital counts must sum to
\code{norb} exactly (the engine refuses otherwise), and the masks override
the \code{nel}/\code{norb} lines --- the menu derives those lines from the
mask, so a mask always describes the space that will actually run.

\section{Walkthrough}

\lstinputlisting[style=fbkcode, firstline=33, lastline=63,
  title={The menu-28 example session (N$_2$, the unconstrained ORMAS mask
  \code{6: 2 0 4, 2 0 4, 2 0 4} on the \%casscf route: the generated input,
  then the measured facts printed with every mask)}]
  {examples/expected/m28_rasormas.txt}

\section{What you get}

\file{<stem>.rasormas.inp} (the \code{\%casscf} route) or
\file{<stem>.rasci.inp} (the CI-only route), plus the run guidance on
screen: check the partition in the output (the \code{\%casscf} route prints
its configuration counts, the \code{\%rasci} route echoes
\code{Number of active orbitals} and \code{First active orbital}); the
orbital optimization omits the active-active rotation for incomplete model
spaces (the manual's own warning, with the \code{actorbs} /
\code{actconstraints} pointers); the \code{\%rasci} module is a standalone
CI that runs no orbital optimization.

\section{Options worth knowing}

\begin{readbox}
\begin{itemize}
  \item \textbf{the route is the physics}: \code{casscf} optimizes the
    orbitals inside the partition (RASSCF / ORMAS-SCF\@: an MCSCF energy);
    \code{rasci} diagonalizes in a fixed orbital window after the frozen
    core (the manual's MRCISD-style usage; \code{ExcLevel} adds excitations
    on top of the references, and \code{mult}/\code{nroots} take the CASSCF
    list syntax);
  \item \textbf{measured anchors} (N$_2$/def2-SVP at 1.10~\AA, fixtures
    \file{n2\_ras.*}, \file{n2\_ormas.*}, \file{n2\_rasci.*}): the
    unconstrained ORMAS mask reproduces the plain CASSCF(6,6) energy to all
    12 printed digits ($-108.989034756374$~Eh) --- so a deviation from the
    CAS energy is exactly the restriction you asked for; the restricted
    \code{RAS(6:2 2/2/2 2)} raises the MCSCF energy to
    $-108.985623236851$~Eh; the CI-only route gives $-108.921051085808$~Eh
    (RAS) and $-108.921178474942$~Eh for the superset ORMAS space;
  \item \textbf{the override guard}: ORCA silently overwrites
    \code{nel}/\code{norb} from either mask (measured for both masks; the
    manual says so for ORMAS); the generator refuses inconsistent numbers
    with that explanation instead, so the input you get always carries the
    mask's own parameters;
  \item \textbf{not wired}: arbitrary-CFG references
    (\code{\{2 2 2 0 0 0\}}), \code{irrep} lists, and the \code{\%rasci}
    module's QDPT/OPA couplings (\code{rel}/\code{douv}) --- all documented
    in the same manual sections and reachable by editing the generated
    file.
\end{itemize}
\end{readbox}
